% Standalone section for problems/3.2/proof.tex — \input this file.
% Written 2026-07-31 (dm window).  Every statement here has a runnable audit in
% problems/3.2/research/scripts/; the audit is named at the end of each item.
% NOTE: this is a NEW file, not an edit of proof.tex, so that the paper keeps a single writer.

\section{Structure of the marked scalar, and the exact location of the obstruction}
\label{sec:marked-scalar-structure}

Throughout, $\Lambda(x,y,z)=\dfrac{(1+x)(1+y)(1+z)\bigl((1+y)(1+z)+xyz\bigr)}{xyz}
=\sum_{\kappa\in P}\lambda_\kappa X^\kappa$ is the Ap\'ery Laurent polynomial,
$b_r=\operatorname{CT}\Lambda^r$, and $C_M(d)$ denotes the $d$-section of the coefficient
array of $\Lambda^M$.

\subsection{The marked scalar is a moment of point counts}

\begin{proposition}\label{prop:moment}
For every prime $p\ge5$ and every $M$,
\[
  C_M(p-1)\;\equiv\;-\sum_{x,y,z\in\mathbb F_p^\times}\Lambda(x,y,z)^M
  \;\equiv\;-\sum_{t\in\mathbb F_p^\times}t^{\,r}N_p(t)\pmod p,
  \qquad r=M\bmod(p-1),
\]
where $N_p(t)=\#\{(x,y,z)\in(\mathbb F_p^\times)^3:\Lambda(x,y,z)=t\}$.
Consequently, writing $n=p+r$, the target condition $p\mid b_n$ is exactly the vanishing
modulo $p$ of the $r$-th moment of the point-count function of the Ap\'ery family.
\end{proposition}

\noindent
\emph{Audit:} \texttt{q32\_marked\_scalar\_character\_sum.py} — $95$ shell-versus-exponential-sum
checks, $95$ shell-versus-moment checks, $44$ shell-versus-Ap\'ery checks.

\subsection{The reflection law, and the parity of $|Z_p|$}

Proposition~\ref{prop:moment} converts a geometric symmetry into an arithmetic one.

\begin{proposition}\label{prop:palindrome}
$N_p(t)=N_p(t^{-1})$ for every $t\in\mathbb F_p^\times$; hence
\[
  b_{\,p-1-r}\equiv b_r\pmod p\qquad(0\le r<p-1),
\]
so $Z_p=\{r<p:\ p\mid b_r\}$ is stable under the involution $r\mapsto p-1-r$, and
$|Z_p|$ is even unless the fixed point $\tfrac{p-1}2$ lies in $Z_p$.
\end{proposition}

Inversion-invariance is not new to the project: it is the hypothesis defining the space
$V_p^+$ in the random model recorded in \texttt{DOCTRINE.md} (Q5387), which yields
$K\sim\mathrm{Bin}((p-3)/2,1/p)\to\mathrm{Poisson}(1/2)$.  What
Proposition~\ref{prop:palindrome} adds is that the invariance is a \emph{theorem} for the
actual sequence, derived from the geometry through Proposition~\ref{prop:moment}, together
with the exact classification of the exceptions: the Poisson law is obeyed by the
\emph{pair} count $|Z_p|/2$ because the involution forces the pairing, and the
``central exceptions'' to evenness are exactly the fixed-point primes.  For $p<4000$ these
are $p=11$ with $Z_{11}=\{5\}$ and $p=3137$ with $Z_{3137}=\{1568\}$, and in both cases the
unique zero equals $(p-1)/2$.

\noindent
\emph{Audit:} same script — $420$ inversion checks, $420$ palindromy checks (the ratio
$b_{p-1-r}/b_r$ is identically $1$, not merely $\pm1$), and the involution verified on all
$548$ primes $5\le p<4000$ without exception.

\subsection{Where $G_n$ actually lives}

The factorisation below is already recorded in \texttt{DOCTRINE.md} (Q5386); what is added
here is the direct computation of both factors.

\begin{proposition}\label{prop:Gn}
Write $a_n=A_n/D_n$ in lowest terms, so $D_n\mid d_n$.  Then
$G_n=\gcd(d_na_n,d_nb_n)=(d_n/D_n)\cdot\gcd(A_n,b_n)$, and for every $n\le330$:
\[
  \gcd(A_n,b_n)=1,
\]
so $G_n=d_n/D_n$ exactly.  Moreover the large prime factors of $G_n$ are precisely the
targets $p\in(n/2,n]$ with $p\mid b_n$; e.g.
\[
  G_{200}=2\cdot3^3\cdot5\cdot17\cdot19\cdot\mathbf{139}\cdot\mathbf{181},\qquad
  G_{321}=2\cdot3\cdot5\cdot7^2\cdot\mathbf{179}\cdot\mathbf{193}\cdot\mathbf{211},
\]
and $G_n$ never exceeds $34$ bits on that range.
\end{proposition}

So the conjecture is equivalent to the statement that Ap\'ery's classical denominator
$d_n=\operatorname{lcm}(1,\dots,n)^3$ is optimal up to $e^{o(n)}$, the only losses coming from
top-window primes dividing $b_n$.

\noindent\emph{Audit:} \texttt{q32\_actual\_Gn\_audit.py}.

\subsection{The target count}

\begin{proposition}\label{prop:K}
Let $K(n)=\#\{p\in(n/2,n]:p\mid b_n\}$, so that the quantity to be bounded is
$T(n)=\sum_{p}\log p\le K(n)\log n$.  Then
\[
  K(n)\le3\qquad\text{for every }n\le200\,000,
\]
with $412$ indices having $K\ge2$, exactly $12$ having $K=3$, and none $K\ge4$; the running
maximum increases only at $n=6,200,321$.  The mean of $K$ is $\approx0.073$, matching the
Mertens prediction $\log2/\log n$.  The same computation gives
$\sum_{p\le200\,000}|Z_p|=18\,126$ (mean $1.016$) and $\max|Z_p|=12$ at $p=159\,977$.
\end{proposition}

\noindent\emph{Audit:} \texttt{q32\_k\_scan.c}, validated against the independent Python
implementation \texttt{q32\_top\_window\_target\_counts.py} on $n\le20\,000$.

\subsection{The denominator-defect law, and what it discharges unconditionally}
\label{subsec:defect-law}

Write $G_n=\prod_p p^{e_p(n)}$ with $e_p(n)=v_p(d_n)-v_p(D_n)\ge0$; since $D_n\mid d_n$ we
always have
\[
  0\le e_p(n)\le v_p(d_n)=3\lfloor\log_p n\rfloor .
\tag{$\ast$}
\]

\begin{proposition}[verified law]\label{prop:defect}
For primes $p\ge7$ with $n/2<p\le n$,
\[
  e_p(n)>0\iff p\mid b_n ,
\]
and then $e_p(n)=1$.  Verified for all such pairs with $n\le250$: $3121$ of $3124$ cases,
the three exceptions being $p=5$ at $n=5,7,9$ — exactly the prime dividing $b_1=5$, which is
the same exception carried by the Ap\'ery--Lucas step.  For $\sqrt n<p\le n/2$ the defect
takes values in $\{1,2,3\}$ and occurs precisely when a base-$p$ digit of $n$ lies in
$Z_p$; e.g.\ the defect residues are $\{5\}$ for $p=11$, $\{3,13\}$ for $p=17$,
$\{8,10\}$ for $p=19$, $\{8,22\}$ for $p=31$, $\{17,19\}$ for $p=37$, $\{10,30\}$ for
$p=41$, $\{9,49\}$ for $p=59$ — in each case exactly $Z_p$, and each pair summing to $p-1$
as Proposition~\ref{prop:palindrome} requires.
\end{proposition}

Two consequences, the first of which needs no hypothesis at all.

\begin{lemma}[unconditional]\label{lem:small-primes}
$\displaystyle\sum_{p\le n/\log n}e_p(n)\log p\;\le\;3\,\pi\!\left(\frac n{\log n}\right)\log n
 \;=\;O\!\left(\frac n{\log n}\right)\;=\;o(n).$
\end{lemma}

\begin{proof}
By $(\ast)$, $e_p(n)\log p\le3\lfloor\log_p n\rfloor\log p\le3\log n$; sum over the
$\pi(n/\log n)\sim n/\log^2n$ primes involved.
\end{proof}

So \emph{every prime below $n/\log n$ is already harmless}, whatever the arithmetic of $b_n$.
For $p>n/\log n$ (hence $p>\sqrt n$ for large $n$) the base-$p$ expansion of $n$ has exactly
two digits, $n=qp+k$ with $1\le q<\log n$, and Ap\'ery--Lucas gives $b_n\equiv b_qb_k\pmod p$.
The digit $q$ can only contribute primes dividing one of $b_1,\dots,b_{\lfloor\log n\rfloor}$,
and
\[
  \sum_{j\le\log n}\log\operatorname{rad}(b_j)\le\sum_{j\le\log n}\log b_j=O(\log^2n),
\]
since $\log b_j\sim3.4j$ (checked: $\log\operatorname{rad}(b_j)$ equals $\log b_j$ to within a
few percent for $j\le14$).  Hence, granting Proposition~\ref{prop:defect},
\[
  \boxed{\;\log G_n\;=\;O\!\left(\frac n{\log n}\right)+O(\log^2n)
   +3\!\!\sum_{\substack{n/\log n<p\le n\\ (n\bmod p)\in Z_p}}\!\!\log p\;}
\]
and the conjecture is \emph{equivalent} to the last sum being $o(n)$.  This is a sharper
localisation than the top-half reduction: the entire range $p\le n/\log n$ and the whole
$q$-digit contribution are discharged with Chebyshev's bound alone.

\subsection{How large the remaining sum actually is}

Using the Lucas criterion ($p\mid b_n$ iff some base-$p$ digit of $n$ lies in $Z_p$) one can
compute the full quantity $R(n)=\log\operatorname{rad}_{p\le n}(b_n)$ for every $n$ in a
range.  For $n\le20\,000$:
\[
\begin{array}{c|c}
 \text{range}&\max R(n)/n\\ \hline
 100\le n\le1000&0.109\ (n=200)\\
 1000\le n\le10\,000&0.0251\ (n=1041)
\end{array}
\qquad
\max_{n\le20\,000}\#\{p\le n:p\mid b_n\}=11 .
\]
Thus $R(n)$ itself grows extremely slowly — the data are consistent with
$R(n)=O(\log n\log\log n)$, which is the heuristic prediction from
$\Pr[p\mid b_n]\approx|Z_p|\log_pn/p$ — while the conjecture only asks for $o(n)$.

\subsection{The obstruction, stated precisely}

The results above localise the difficulty completely: $G_n$ is a handful of small primes
times the $K(n)\le3$ targets, and every route we have tried terminates in the same place —
bounding the number of prime factors of the \emph{single} integer $b_n$ inside the short
window $(n/2,n]$.  Three separate barriers are now identified, each with a reason.

\begin{enumerate}
\item \textbf{Rank one is forced, not accidental.}  The Newton polytope of $\Lambda$ has a
unique interior lattice point (the origin), so by Beukers--Vlasenko \emph{Dwork Crystals I}
(IMRN 2021, Thm.~11) the Cartier-stable quotient has rank $1$; equivalently, for the K3
fibre $h^{0,2}=1$, so the Hasse--Witt operator is $1\times1$ (Katz, \emph{Une formule de
congruence pour la fonction z\^eta}, SGA 7 II, Exp.~XXII).  Hence \emph{every} mod-$p$
observable built from the coefficient array is a multiple of the single marked scalar of
Proposition~\ref{prop:moment}, and any elimination of it reacquires the candidate primorial,
of logarithmic height $\sim n/2$.  This is the structural reason behind every rank-one alias
encountered in the terminal, Turán, Hankel, Newton-jet and cross-index families.

\item \textbf{The second $p$-adic digit is not a second selective condition.}  Modulo $p^2$
one does obtain a genuinely new coordinate, namely $D_{p,r}=R_{p-1}(\Lambda^{r-1}\Phi_p)$
with $\Phi_p=(\Lambda^p-\Lambda(X^p))/p$; but it is not target-selective.  Direct computation
gives $D_{19,10}=D_{31,8}=0$ at genuine targets, and also $D_{19,9}=0$ at a non-target.  No
codimension-two statement is available.

\item \textbf{The vertical average is equally out of reach.}  The pointwise problem would
follow from Hypothesis~Z, $\sum_{p\le P}|Z_p|=O(P/\log P)$.  The best unconditional aggregate
implied by $|Z_p|=O(p^{2/3})$ is $O(P^{5/3}/\log P)$, and no improvement is available — not to
$P^{1+\varepsilon}/\log P$, and not even $|Z_p|=O(1)$ on a positive-density set of primes.
Katz's finite-field Mellin equidistribution (\emph{Convolution and Equidistribution},
Ann.\ of Math.\ Studies 180, 2012) controls archimedean distributions of normalised traces,
not exact vanishing modulo $p$; and Serre's bound
$\#\{p\le X:a_p(E)=0\}\ll X/(\log X)^{3/2-\varepsilon}$
(Publ.\ Math.\ IHES 54 (1981)) is one trace test per prime, whereas $|Z_p|$ involves $p$
character-indexed tests inside each prime.
\end{enumerate}

Accordingly the unconditional pointwise statement appears to require genuinely new
technology — a characteristic-varying local-limit or large-sieve theorem for the
two-variable Mellin sheaf — rather than a refinement of the arguments used here.
